\subsection*{Study B: on held-out analyses both review styles find every documented error, and both over-call errors in correct analyses}
In 132 reviews of the 22 held-out items, every review of a flawed item
found its documented error, in both arms. This held for all 11 flawed
items and for 33 of 33 reviews per arm (Wilson 95\% interval over items
74--100\%; items are the unit because the three reviews of an item are
not independent). It held whether the error fell inside the
checklist's patterns (6 items) or outside them (5 items: a circular
correlation that fails on evenly spread angles, span-filled residue
labels, an arbitrary sign flip of singular vectors, raw counts fed to an
encoder that expects binned values, and a row-index join that matched the
wrong cells). Both graders judged all 66 reviews of flawed items to have
found the documented error, so they agreed on all 66. Cohen's $\kappa$
is not defined when every judgment is the same, so we report this raw
agreement instead.

On the 11 correct (clean) items, most reviews still claimed at least
one error in a main result that the item's key does not count as an
error (Fig~\ref{fig:studyB}A). This happened in 79\% of checklist
reviews (26 of 33; 95\% bootstrap interval over items 58--94\%) and in
82\% of generic reviews (27 of 33; 61--100\%). The paired difference,
checklist minus generic, was $-3$ percentage points (95\% bootstrap
interval over items $-21$ to $+15$; Wilcoxon $p=0.88$). Checklist
reviews made 45 such claims in all (1.4 per review of a clean item), and
generic reviews made 43 (1.3 per review). Checklist reviews were longer
(9.2 against 7.8 findings per review); both arms found about two real
problems per review beyond the documented error (2.1 in each).

\begin{figure}[H]
\caption{\textbf{Study B: held-out analyses from other projects.} Every
review detected the documented error in every flawed item, in both arms
(33 of 33 each; not plotted). (A)~For each of the 11 correct (clean)
items, the number of its three reviews per arm that claimed at least one
error in a main result that the item's key does not count as an error.
(B)~Post-hoc classification of the 82 such claims that the graders
marked as false alarms, on clean and flawed items, by arm. The 23
further claims that the verifier ruled false alarms after a grader
dispute are not included.}
\label{fig:studyB}
\end{figure}

Across clean and flawed items, the reviews made 105 such claims (88 on
clean items, 17 on flawed items). The graders marked 82 of them as false
alarms (73 on clean items, 9 on flawed items). The other 23 (15 on clean
items, 8 on flawed items) were findings that a grader marked as disputed
and the verifier agent then ruled false alarms. After the results were
known, we classified the 82 claims that the graders marked. This step was
not pre-registered. One agent labelled each claim, with no second rater.
Most of the 82 were accurate observations whose severity was overstated
(Fig~\ref{fig:studyB}B): 59 were correct but did not change any stated
conclusion (for example, that an input-label overlap weakened, but did
not overturn, an interpretation), 17 were correct and already stated in
the analysis's own summary or README (for example, that a replication in
a second cell line rested mostly on one transcription factor), 3 were
factually wrong and 3 were real problems that the key had missed. The
pattern was the same in both arms. The 23 claims ruled by the verifier
were not classified.

